% Options for packages loaded elsewhere
\PassOptionsToPackage{unicode}{hyperref}
\PassOptionsToPackage{hyphens}{url}
\documentclass[
]{article}
\usepackage{xcolor}
\usepackage{amsmath,amssymb}
\setcounter{secnumdepth}{-\maxdimen} % remove section numbering
\usepackage{iftex}
\ifPDFTeX
  \usepackage[T1]{fontenc}
  \usepackage[utf8]{inputenc}
  \usepackage{textcomp} % provide euro and other symbols
\else % if luatex or xetex
  \usepackage{unicode-math} % this also loads fontspec
  \defaultfontfeatures{Scale=MatchLowercase}
  \defaultfontfeatures[\rmfamily]{Ligatures=TeX,Scale=1}
\fi
\usepackage{lmodern}
\ifPDFTeX\else
  % xetex/luatex font selection
\fi
% Use upquote if available, for straight quotes in verbatim environments
\IfFileExists{upquote.sty}{\usepackage{upquote}}{}
\IfFileExists{microtype.sty}{% use microtype if available
  \usepackage[]{microtype}
  \UseMicrotypeSet[protrusion]{basicmath} % disable protrusion for tt fonts
}{}
\makeatletter
\@ifundefined{KOMAClassName}{% if non-KOMA class
  \IfFileExists{parskip.sty}{%
    \usepackage{parskip}
  }{% else
    \setlength{\parindent}{0pt}
    \setlength{\parskip}{6pt plus 2pt minus 1pt}}
}{% if KOMA class
  \KOMAoptions{parskip=half}}
\makeatother
\usepackage{longtable,booktabs,array}
\newcounter{none} % for unnumbered tables
\usepackage{multirow}
\usepackage{calc} % for calculating minipage widths
% Correct order of tables after \paragraph or \subparagraph
\usepackage{etoolbox}
\makeatletter
\patchcmd\longtable{\par}{\if@noskipsec\mbox{}\fi\par}{}{}
\makeatother
% Allow footnotes in longtable head/foot
\IfFileExists{footnotehyper.sty}{\usepackage{footnotehyper}}{\usepackage{footnote}}
\makesavenoteenv{longtable}
\usepackage{graphicx}
\makeatletter
\newsavebox\pandoc@box
\newcommand*\pandocbounded[1]{% scales image to fit in text height/width
  \sbox\pandoc@box{#1}%
  \Gscale@div\@tempa{\textheight}{\dimexpr\ht\pandoc@box+\dp\pandoc@box\relax}%
  \Gscale@div\@tempb{\linewidth}{\wd\pandoc@box}%
  \ifdim\@tempb\p@<\@tempa\p@\let\@tempa\@tempb\fi% select the smaller of both
  \ifdim\@tempa\p@<\p@\scalebox{\@tempa}{\usebox\pandoc@box}%
  \else\usebox{\pandoc@box}%
  \fi%
}
% Set default figure placement to htbp
\def\fps@figure{htbp}
\makeatother
\ifLuaTeX
  \usepackage{luacolor}
  \usepackage[soul]{lua-ul}
\else
  \usepackage{soul}
\fi
\setlength{\emergencystretch}{3em} % prevent overfull lines
\providecommand{\tightlist}{%
  \setlength{\itemsep}{0pt}\setlength{\parskip}{0pt}}
\ifPDFTeX
  \TeXXeTstate=1
  \newcommand{\RL}[1]{\beginR #1\endR}
  \newcommand{\LR}[1]{\beginL #1\endL}
  \newenvironment{RTL}{\beginR}{\endR}
  \newenvironment{LTR}{\beginL}{\endL}
\fi
\ifluatex
  \newcommand{\RL}[1]{\bgroup\textdir TRT#1\egroup}
  \newcommand{\LR}[1]{\bgroup\textdir TLT#1\egroup}
  \newenvironment{RTL}{\textdir TRT\pardir TRT\bodydir TRT}{}
  \newenvironment{LTR}{\textdir TLT\pardir TLT\bodydir TLT}{}
\fi
\usepackage{bookmark}
\IfFileExists{xurl.sty}{\usepackage{xurl}}{} % add URL line breaks if available
\urlstyle{same}
% fallback for those not using the hyperref driver hyperxmp:
\makeatletter
\@ifundefined{xmpquote}{\newcommand{\xmpquote}[1]{#1}}{}
\makeatother
\hypersetup{
  hidelinks,
  pdfcreator={LaTeX via pandoc}}

\author{}
\date{}

\begin{document}


\includegraphics[width=0.98194in,height=0.27014in]{media/image1.jpeg}J.
Stat. Appl. Pro. Vol. No. (20-\/-) 185

\begin{quote}
Journal of Statistics Applications \& Probability

\emph{An International Journal}

http://dx.doi.org/10.18576/jsap/1001-\/-
\end{quote}

\textbf{Statistical Analysis of Body Dysmorphic Disorder, Quality of
Life, and Social Anxiety in Fire Victims: Path Analysis}

Abdallah Almahaireh\textsuperscript{1}, Baha\textquotesingle{}
Shawaqfeh\textsuperscript{2}, and Bahaa Eldeen Abu
Serdaneh\textsuperscript{3}

\textsuperscript{1} Professor in Counselling and Special Education
Department, School of Educational Science, The University of Jordan.

\textsuperscript{2} Full-time lecturer in Counselling and Mental Health
Department, Faculty of Educational Science, The World Islamic Sciences
and Education University

\textsuperscript{3} Ph.D. in Educational Psychology, School of
Educational Science, The University of Jordan.

Received: 25 Nov. 2022, Revised: 4 Jan. 2023, Accepted: 15 Jan. 2023.

Published online: 1 -\/-\/-. 20-\/-.

\textbf{Abstract:} The study aimed to explore the relationship between
body dysmorphic disorder symptoms and quality of life, and the role of
social anxiety disorder symptoms among fire victims. The study used a
quantitative technique using path analysis. The study sample consists of
292 participants. The study used body dysmorphic disorder symptoms,
quality of life, and social anxiety disorder symptoms scales. The
results showed that there is a moderate level of quality of life among
fire victims while there is a High level of body dysmorphic disorder
symptoms and social anxiety symptoms. Also, the results show that social
anxiety disorder symptoms have moderate role in the relationship between
body dysmorphic disorder symptoms and quality of life among fire
victims. On the other hand, body dysmorphic disorder symptoms influence
quality of life through one path which is indirect mediated by social
anxiety disorder symptoms. The model explains 46\% of quality of life.
The study concludes that this result should be considered when
specialists need to improve the quality of life among fire victims by
addressing the negative body dysmorphic and focusing on their social
anxiety symptoms. This may help them improve their mental health and
wellness.

\textbf{Keywords:} Body Dysmorphic, Quality of Life, Social Anxiety,
Fire Victims.

\section{1 Introduction}\label{introduction}

Fires is one of the common causes for injuries and death globally, which
occur 180 thousand annually\textsuperscript{{[}1{]}}. In Jordan 21999
patients visits burn unit in governmental hospital
\textsuperscript{{[}2{]}}. There is a lack of interest in burning
injuries which need a long-term treatment which may happen to anyone,
anywhere, and anytime. Burning May be caused by many factors such as
gas, electricity, and hot water which result psychological and physical
problems resulting decreasing in Quality of life (QOL) among fire
victims.

QOL is affected by many variables related to burning such as; the
response to injury, speed treatment response, quality of treatment,
social support, and counselling \textsuperscript{{[}3{]}}. The QOL is
defined as ``the individual's perception of their position in life in
the context of culture and value systems in which they live and in
relation to their goals, expectations, standards, and concerns''
\textsuperscript{{[}4{]}}.

QOL is an extent of individual's awareness of the experiences and how
much he's satisfied with his basic need fulfillment in life. Also, how
much they feel happiness, accomplishment, and integration in terms of
psychological, physical, cognitive, social, and emotional health
\textsuperscript{{[}5{]}}. Fire victim's QOL is affected by many
aspects. Gautam et al., \textsuperscript{{[}6{]}} study found that burn
injuries has a negative impact on QOL and its domains; psychology,
physically, socially, and environmentally.

Fires harm the body physically and leads to marks and deformities which
affect the individual's body image. Fire victims could be experience
body dysmorphic disorder symptoms (BDDS) that is defined as a disorder
caused by a body defect in the external appearance that can be seen by
others or can't be seen\textsuperscript{{[}7{]}}. How the individual
sees him self is very important because it develops his perception of
self and how much he's acceptable to others \textsuperscript{{[}8{]}},
the burning marks can lead the person to be not satisfied of his body
and they know that others can see it which leads him to avoid social
situations \textsuperscript{{[}9{]}}.

The social anxiety disorder symptoms (SADS) refers to a high level of
anxiety happened when the individual faces social engagement and it
leads him to avoid any interaction with others also it effects his life
socially, occupationally, or on other functional areas
\textsuperscript{{[}10{]}}. This can be affected negatively on their
quality of life and leads them to isolation. Especially when they review
the negative situations they have been exposed to, then evaluate them
and create a negative thought about it, this affect their professional
and social performance. It's an important matter for the individual to
feel healthy and psychologically wellbeing.

The best way to investigates these relations is by using path analysis
method which is a statistical method investigate the direct, indirect,
and total effects by mediating variable between at least two variables
and it's an extension method of multiple regression
\textsuperscript{{[}11{]}}. This method helps in to test the theoretical
model which assumed in this study that the SADS mediating the
relationship between BDDS and QOL.

This study focused on how the BDDS can affect QOL, and if the SADS has a
moderate role between them among fire victims. We believe that because
of the body deformities that shows on fire victims it may increase BDDS
this leads them to avoid social interactions and expressed SADS that
affects negatively on many functional areas related to the QOL. The
suggested model shown in figure 1.

\emph{Figure 1: The suggested model for relationship between BDDS and
QOL and the moderating of SADS}

\textbf{1.1 Research questions}

What is the level of the BDDS, QOL, and SADS among fire victims?

Is there a relationship between BDDS, QOL, and SADS among fire victims?

Do the SADS moderate the relationship between BDDS and QOL among fire
victims?

\textbf{\hfill\break
}

\begin{enumerate}
\def\labelenumi{\arabic{enumi}.}
\setcounter{enumi}{1}
\item
  \textbf{Materials and Methods}
\end{enumerate}

\emph{2.1 Study design}

This research employs a quantitative technique using path analyses. A
volunteer selected 305 fire victims those who survived the burning from
the Dermatology Clinics in Jordan from 2020 to 2025. The inclusion
criteria for this study were: (1) accepted to by in this research, (2)
exposure to the burning between 2020 and 2025, (3) the burning is effect
on main parts like face, nick, and hands who usually don't cover them,
(4) who shows that the burning affect their life. as subjects; all
participants supplied written informed permission with the scales. The
research ethics committee authorized the institution's study.

\emph{2.2 Participants}

The study sample consists of 292 fire victims from the Dermatology
Clinic in Jordan from 2020 to 2025. The distribution of participants by
socioeconomic rank is as follows: 100 Male and 190 females. 137 were
single and 153 were married. Their age was between 18-50 with a mean of
(33.6 ±8.9).

\emph{2.3 Data collection Tools}

\emph{\textbf{2.3.1 Body Dysmorphic Disorder Symptoms Scale}}

The researchers used a scale of \RL{24} items from Al-Hasi
\textsuperscript{{[}12{]}}. Participants rated each item on a
four-Likert scale, from 1 (Strongly disagree) to 4 (Strongly agree). The
researcher extracted the validity and reliability and found that the
discriminate evidence ranged between 0.344 and 0.713. Cronbach's alpha
was 0.914.

\emph{\textbf{2.3.2 Social Anxiety Disorder Symptoms Scale}}

The researchers used a scale of 10 items from Shawaqfeh \& Almahaireh
\textsuperscript{{[}13{]}}. Participants rated each item on a
four-Likert scale, from 1 (Strongly disagree) to 4 (Strongly agree). The
researcher extracted the validity and reliability and found that the
discriminate evidence ranged between 0.478 and 0.915. Cronbach's alpha
was 0.916.

\emph{\textbf{2.3.3 Quality of Life Scale}}

The researchers used a scale of 24 items from Abu Al-wafaa
\textsuperscript{{[}14{]}}. It's including 4 dimensions (Physical,
Psychological, Social Relationships, Environmental). Participants rated
each item on a four-Likert scale, from 1 (Never) to 4 (Always). The
researcher extracted the validity and reliability and found that the
discriminate evidence ranged between 0.345 and 0.922. Cronbach's alpha
was ranged between 0.737-0.918.

\emph{2.4 Data collection}

The researchers obtained approval for this study from the Institutional
Review Board at Al-Ahliyya Amman University. The registration number
467/2025/1. The date was 17/6/2025. Anonymized data was collected by the
researchers, the survey was conducted by paper instruments, the first
page contain the purpose and the significance of the study.

\emph{2.5 Data analysis}

The researchers entered the data into SPSS data sheet, and used IBM SPSS
Statistics for Windows, version 27 to analyze the data. Means, standard
deviations, Pearson's correlation were calculated. Also, IBM SPSS AMOS
v24 was used to do the path analyses, the missing data were filled using
SPSS program using series mean method in replace missing values, this
procedure used because AMOS don't calculate any data with missing
values.

Path analysis used to examine and explore the relationships between the
observed variables (BDDS, SADS, QOL). This technique is used to test the
link between two variables via a third variable. Because of that it
considered as a complex technique than multiple regression analysis. The
advantage of sing it is its superior ability to detect relationships and
it studied many variables at the same time \textsuperscript{{[}15{]}}.

\begin{enumerate}
\def\labelenumi{\arabic{enumi}.}
\setcounter{enumi}{2}
\item
  \textbf{Findings}
\end{enumerate}

\emph{3.1 The level of the BDDS, QOL, and SADS}

To investigate the level of the BDDS, QOL, and SADS among fire victims.
Means and standard deviations were calculated. Table 1 show that the
level of the QOL among fire victims was moderate with a mean of
(2.29±0.169). While the levels of BDDS and SADS was high with a means of
(3.14±0.568) for the BDDS and (3.08±0.704) for the SADS.

{\def\LTcaptype{none} % do not increment counter
\begin{longtable}[]{@{}
  >{\centering\arraybackslash}p{(\linewidth - 6\tabcolsep) * \real{0.4355}}
  >{\centering\arraybackslash}p{(\linewidth - 6\tabcolsep) * \real{0.1572}}
  >{\centering\arraybackslash}p{(\linewidth - 6\tabcolsep) * \real{0.1862}}
  >{\centering\arraybackslash}p{(\linewidth - 6\tabcolsep) * \real{0.2210}}@{}}
\toprule\noalign{}
\begin{minipage}[b]{\linewidth}\centering
\textbf{Variable}
\end{minipage} & \begin{minipage}[b]{\linewidth}\centering
\textbf{Mean}
\end{minipage} & \begin{minipage}[b]{\linewidth}\centering
\textbf{Std.}
\end{minipage} & \begin{minipage}[b]{\linewidth}\centering
\textbf{Level}
\end{minipage} \\
\midrule\noalign{}
\endhead
\bottomrule\noalign{}
\endlastfoot
\textbf{BDDS} & 3.14 & .568 & High \\
\textbf{QOL} & 2.29 & .169 & Moderate \\
\textbf{SADS} & 3.08 & .704 & High \\
\end{longtable}
}

\emph{Table 1: Means and standard deviations of the BDDS, QOL, and SADS}

\emph{3.2 The relationship between BDDS, QOL, and SADS}

To investigate the relationship between BDDS, QOL, and SADS among fire
victims. Person correlation coefficients were calculated. Table 2 that
there is statistically significant negative relationship between QOL and
each of BDDS and SADS (R= -.626 and -0.644). While there is
statistically significant positive relationship between BDDS and SADS
(R= .675).

{\def\LTcaptype{none} % do not increment counter
\begin{longtable}[]{@{}
  >{\centering\arraybackslash}p{(\linewidth - 12\tabcolsep) * \real{0.2495}}
  >{\centering\arraybackslash}p{(\linewidth - 12\tabcolsep) * \real{0.1321}}
  >{\centering\arraybackslash}p{(\linewidth - 12\tabcolsep) * \real{0.1188}}
  >{\centering\arraybackslash}p{(\linewidth - 12\tabcolsep) * \real{0.1320}}
  >{\centering\arraybackslash}p{(\linewidth - 12\tabcolsep) * \real{0.1178}}
  >{\centering\arraybackslash}p{(\linewidth - 12\tabcolsep) * \real{0.1320}}
  >{\centering\arraybackslash}p{(\linewidth - 12\tabcolsep) * \real{0.1178}}@{}}
\toprule\noalign{}
\multirow{2}{=}{\centering\arraybackslash \begin{minipage}[b]{\linewidth}\centering
\textbf{Variable}
\end{minipage}} & \begin{minipage}[b]{\linewidth}\centering
\textbf{BDDS}
\end{minipage} & \begin{minipage}[b]{\linewidth}\centering
\end{minipage} & \begin{minipage}[b]{\linewidth}\centering
\textbf{QOL}
\end{minipage} & \begin{minipage}[b]{\linewidth}\centering
\end{minipage} & \begin{minipage}[b]{\linewidth}\centering
\textbf{SADS}
\end{minipage} & \begin{minipage}[b]{\linewidth}\centering
\end{minipage} \\
& \begin{minipage}[b]{\linewidth}\centering
\textbf{Correlation}
\end{minipage} & \begin{minipage}[b]{\linewidth}\centering
\textbf{Sig.}
\end{minipage} & \begin{minipage}[b]{\linewidth}\centering
\textbf{Correlation}
\end{minipage} & \begin{minipage}[b]{\linewidth}\centering
\textbf{Sig.}
\end{minipage} & \begin{minipage}[b]{\linewidth}\centering
\textbf{Correlation}
\end{minipage} & \begin{minipage}[b]{\linewidth}\centering
\textbf{Sig.}
\end{minipage} \\
\midrule\noalign{}
\endhead
\bottomrule\noalign{}
\endlastfoot
\textbf{BDDS} & - & - & -.626* & .000 & .675* & .000 \\
\textbf{QOL} & - & - & - & - & -.644* & .000 \\
\textbf{SADS} & - & - & - & - & - & - \\
\end{longtable}
}

\textbf{** sig at the 0.01 level (2-talied)}

\emph{Table 2: Person correlation test results between BDDS, QOL, and
SADS}

\emph{3.3 Examination of the moderating effect of SADS}

To examine if SADS has a moderate effect in the relationship between
BDDS and QOL among fire victims, path analyses were computed using SPSS
AMOS V24. The assumptions were checked by checking the normality
distribution. The results in table 3 and the figures 2-4 show that the
variables are normally distributed which the Skewness ranged between
.398 and -.945, while the Kurtosis ranged between -.745 and .954
\textsuperscript{{[}16{]}}.

{\def\LTcaptype{none} % do not increment counter
\begin{longtable}[]{@{}
  >{\centering\arraybackslash}p{(\linewidth - 8\tabcolsep) * \real{0.3343}}
  >{\centering\arraybackslash}p{(\linewidth - 8\tabcolsep) * \real{0.1768}}
  >{\centering\arraybackslash}p{(\linewidth - 8\tabcolsep) * \real{0.1632}}
  >{\centering\arraybackslash}p{(\linewidth - 8\tabcolsep) * \real{0.1630}}
  >{\centering\arraybackslash}p{(\linewidth - 8\tabcolsep) * \real{0.1626}}@{}}
\toprule\noalign{}
\multirow{2}{=}{\centering\arraybackslash \begin{minipage}[b]{\linewidth}\centering
\textbf{Variable}
\end{minipage}} &
\multicolumn{2}{>{\centering\arraybackslash}p{(\linewidth - 8\tabcolsep) * \real{0.3400} + 2\tabcolsep}}{%
\begin{minipage}[b]{\linewidth}\centering
\textbf{Skewness}
\end{minipage}} &
\multicolumn{2}{>{\centering\arraybackslash}p{(\linewidth - 8\tabcolsep) * \real{0.3257} + 2\tabcolsep}@{}}{%
\begin{minipage}[b]{\linewidth}\centering
\textbf{Kurtosis}
\end{minipage}} \\
& \begin{minipage}[b]{\linewidth}\centering
\textbf{Statistic}
\end{minipage} & \begin{minipage}[b]{\linewidth}\centering
\textbf{Std. Error}
\end{minipage} & \begin{minipage}[b]{\linewidth}\centering
\textbf{Statistic}
\end{minipage} & \begin{minipage}[b]{\linewidth}\centering
\textbf{Std. Error}
\end{minipage} \\
\midrule\noalign{}
\endhead
\bottomrule\noalign{}
\endlastfoot
\textbf{BDDS} & -.945 & .143 & .789 & .285 \\
\textbf{QOL} & .398 & .143 & .954 & .285 \\
\textbf{SADS} & -.426 & .143 & -.745 & .285 \\
\end{longtable}
}

\emph{Table 3: Means and standard deviations of the BDDS, QOL, and SADS}

\includegraphics[width=3.95586in,height=2.32576in]{media/image5.png}

\emph{Figure 2: The normally distribution of the QOL}

It's appears from figure 2 that the QOL is normally distributed which
the histogram show that the values ranged between 1 to 3. While the most
values ranged between 2.1 to 2.6 and it's appears that there are normal
distribution.

\includegraphics[width=3.99242in,height=2.34726in]{media/image6.png}

\emph{Figure 3: The normally distribution of the BDDS}

It's appears from figure 3 that the BDDS is normally distributed which
the histogram show that the values ranged between 1.4 to 4. While the
most values ranged between 2.4 to 3.9 and it's appears that there are
normal distribution.

\includegraphics[width=3.87879in,height=2.28045in]{media/image7.png}

\emph{Figure 4: The normally distribution of the SADS}

It's appears from figure 4 that the SADS is normally distributed which
the histogram show that the values ranged between 1 to 3. While the most
values ranged between 1.7 to 4 and it's appears that there is normal
distribution.

Also, the model fit parameters were checked in table 4. The results show
that the model is fit due to the CMIN/df=2.428, RMSEA= 0.069, and CFI=
0.739 according to Hair et al. \textsuperscript{{[}16{]}}.

{\def\LTcaptype{none} % do not increment counter
\begin{longtable}[]{@{}
  >{\centering\arraybackslash}p{(\linewidth - 4\tabcolsep) * \real{0.2931}}
  >{\centering\arraybackslash}p{(\linewidth - 4\tabcolsep) * \real{0.2859}}
  >{\centering\arraybackslash}p{(\linewidth - 4\tabcolsep) * \real{0.2926}}@{}}
\toprule\noalign{}
\begin{minipage}[b]{\linewidth}\centering
\textbf{Model fit parameters}
\end{minipage} & \begin{minipage}[b]{\linewidth}\centering
\textbf{Value}
\end{minipage} & \begin{minipage}[b]{\linewidth}\centering
\textbf{Acceptable value}
\end{minipage} \\
\midrule\noalign{}
\endhead
\bottomrule\noalign{}
\endlastfoot
\textbf{CMIN/df} & 2.428 & Less than 5 \\
\textbf{RMSEA} & 0.069 & Less than 0.1 \\
\textbf{CFI} & 0.739 & Greater than 0.7 \\
\end{longtable}
}

\emph{Table 4: Model fit parameters}

In figure 5 the results show that SADS has moderate role in the
relationship between BDDS and QOL among fire victims. According to table
5 BDDS influence QOL through one path which is indirect mediated by
SADS. Although, there are positive direct effect of BDDS on SADS
achieves .459, and negative direct effect of SADS on QOL achieves -.654.
On the other hand, the direct effect of BDDS on QOL was quietly low
(-.049) but the indirect effect was significant and negative which was
-.300.

\includegraphics[width=4.02778in,height=5.09325in]{media/image8.png}

\emph{Figure 5: The tested model for relationship between BDDS and QOL
and the moderating of SADS}

{\def\LTcaptype{none} % do not increment counter
\begin{longtable}[]{@{}
  >{\centering\arraybackslash}p{(\linewidth - 8\tabcolsep) * \real{0.1892}}
  >{\centering\arraybackslash}p{(\linewidth - 8\tabcolsep) * \real{0.1892}}
  >{\centering\arraybackslash}p{(\linewidth - 8\tabcolsep) * \real{0.1892}}
  >{\centering\arraybackslash}p{(\linewidth - 8\tabcolsep) * \real{0.1892}}
  >{\centering\arraybackslash}p{(\linewidth - 8\tabcolsep) * \real{0.1892}}@{}}
\toprule\noalign{}
\begin{minipage}[b]{\linewidth}\centering
\textbf{Independent}
\end{minipage} & \begin{minipage}[b]{\linewidth}\centering
\textbf{Dependent}
\end{minipage} & \begin{minipage}[b]{\linewidth}\centering
\textbf{Direct effect}
\end{minipage} & \begin{minipage}[b]{\linewidth}\centering
\textbf{Indirect effect}
\end{minipage} & \begin{minipage}[b]{\linewidth}\centering
\textbf{Total effect}
\end{minipage} \\
\midrule\noalign{}
\endhead
\bottomrule\noalign{}
\endlastfoot
\textbf{BDDS} & \textbf{SADS} & .459 & - & .459 \\
\textbf{BDDS} & \textbf{QOL} & -.049 & -.300 & -.349 \\
\textbf{SADS} & \textbf{QOL} & -.654 & - & -.654 \\
\end{longtable}
}

\emph{Table 5: Direct, indirect, and Total effects between the
variables}

After the results showed that SADS has moderate role in the relationship
between BDDS and QOL among fire victims. The square multiple regression
in the model was investigate. As shown in table 6 the model explains
21.1\% from SADS and 46\% from QOL. It is a good weight of prediction.

{\def\LTcaptype{none} % do not increment counter
\begin{longtable}[]{@{}
  >{\centering\arraybackslash}p{(\linewidth - 2\tabcolsep) * \real{0.4717}}
  >{\centering\arraybackslash}p{(\linewidth - 2\tabcolsep) * \real{0.4717}}@{}}
\toprule\noalign{}
\begin{minipage}[b]{\linewidth}\centering
\textbf{Variable}
\end{minipage} & \begin{minipage}[b]{\linewidth}\centering
\textbf{Estimate}
\end{minipage} \\
\midrule\noalign{}
\endhead
\bottomrule\noalign{}
\endlastfoot
\textbf{SADS} & .211 \\
\textbf{QOL} & .460 \\
\end{longtable}
}

\emph{Table 6: square multiple regression for SADS and QOL in the model}

\begin{enumerate}
\def\labelenumi{\arabic{enumi}.}
\setcounter{enumi}{3}
\item
  \textbf{Discussion}
\end{enumerate}

The results of this study show that the SADS has a moderate role in the
relationship between BDDS and QOL among fire victims. Thus, the BDDS has
only indirect effect through the SADS but the direct effect in the model
was very low despite that the correlation between them was significant
and negative -.626. On the other hand, the fire victims have a moderate
level of QOL and High level of SADS and BDDS.

These results show that the burn increases the BDDS, makes the person
think others will not accept them, and look at them negatively. they
focus on how others will judge them, and they feel humiliated when
others look at them because of their burn marks
\textsuperscript{{[}17{]}}. SADS makes QOL decrease because it affects
how they deals with others and tries to avoid them, leading them to stay
alone \textsuperscript{{[}18{]}}. Also, they may avoid going to their
workplace because they feel ashamed and fears how others will assess
them and their burn marks \textsuperscript{{[}19{]}}, which affect their
occupational and social QOL. Emotionally, they suffered from negative
feelings, pessimism, and low self-esteem related to the negative
thinking about their body \textsuperscript{{[}20{]}}.

Usually, fire victims struggled with others, such as when others look at
them intently or appear afraid. This means that self-perception is the
most important factor in how BDDS affect QOL. If the fire victims feel
afraid from others staring at them and think that their burns are a
stigma, they will suffer from SADS, which leads them to avoidance and
depression overall; this will decrease QOL.

The Strengths of this study are its the first research study the
relations BDDS, QOL, SADS. Used the path analyses and structural
equation model to investigate the SADS as a moderate variable between
BDDS and QOL. The results was very impressive in this study especially
that when SADS moderate the relationship between BDDS and QOL it
canceled the direct impact and change it to indirect impact even though
the relation was -.626 which consider high correlation value.

We believe that this study can be improved by entering other variables
in the model such as self-talk between BDDS and SADS, or self-perception
which we think that it has a good impact on QOL. Also, the researcher
should consider compare the fire victims with other participants to make
sure if these results applied on them or all how had Body deformity and
focus on compare between the individuals who had Body deformity from
birth and who's involved in an accident.

\begin{enumerate}
\def\labelenumi{\arabic{enumi}.}
\setcounter{enumi}{4}
\item
  \textbf{Conclusions}
\end{enumerate}

In conclusion, we believe that the model has proved a good explanation
for the effect of BDDS on QOL through SADS. It should be considered in
future research and when specialists need to improve the quality of life
among fire victims by addressing the negative body dysmorphic and
focusing on their social anxiety symptoms. This may help them improve
their mental health and wellness. Further studies could investigate the
self-esteem, self-perception, and social self-efficacy.

\textbf{Data Availability:}

The datasets analyzed during the current study are available from the
corresponding author upon reasonable request.

\textbf{Conflict of Interest:}

The authors declare no conflicts of interest in this study.

\textbf{Funding:}

This research received no specific grants from funding agencies.

\textbf{Consent to Participate:}

Participants provided informed consent to participate in this study.

\textbf{Ethics approval:}

This study followed the principles of the Committee on Publication
Ethics (COPE).

\textbf{\hfill\break
}

\textbf{References}

\begin{enumerate}
\def\labelenumi{\arabic{enumi}.}
\item
  World Health Organization (WHO) (2023). \emph{Burns}. Available from:
  \href{https://www.who.int/news-room/fact-sheets/detail/burns}{\ul{https://www.who.int/news-room/fact-sheets/detail/burns}}
\item
  Jordanian Ministry of Health (MOH) (2025). \emph{Annual statistically
  report for 2024}.
\item
  Hunt JL, Arnoldo BD, Purdue GF. Chapter 4---prevention of burn
  injuries. \emph{Total burn care}. 7th ed. London: WB Saunders.
  2012:47‑55.
\item
  Ho, J. H. (2024). \emph{University of Washington quality of life
  questionnaire}. In Encyclopedia of Quality of Life and Well-Being
  Research (pp. 7398-7400). Cham: Springer International Publishing.
\item
  Al-Harthi, I. \& Al-Samadi, A. (2017). The level of health behavior
  among Umm Al-Qura University students: A descriptive study,
  \emph{Educational Journal}, 31 (122), 125-145.
\item
  Gautam, R., Rajoura, O. P., Sharma, A. K., \& Bhatia, M. S. (2022).
  Socio-demographic features and quality of life post burn injury.
  \emph{Journal of family medicine and primary care}, 11(3), 1032-1035.
\item
  Rück, C., Mataix-Cols, D., Feusner, J. D., Shavitt, R. G., Veale, D.,
  Krebs, G., \& Fernández de la Cruz, L. (2024). Body dysmorphic
  disorder. \emph{Nature Reviews Disease Primers}, 10(1), 92.
\item
  Sreshta, N., Pope, H. G., Hudson, J. I., \& Kanayama, G. (2017).
  Muscle dysmorphia. \emph{Body dysmorphic disorder: advances in
  research and clinical practice}, 81.
\item
  Alhammadi, A. (2020). \emph{Overview of the Diagnostic and Statistical
  Manual of Mental Disorders}, Alarabia for science.
\item
  American Psychiatric Association. (2013). \emph{Diagnostic and
  statistical manual of mental disorders} (DSM-5®). American Psychiatric
  Pub.
\item
  Columbia (2026). Path Analysis. Available from:
  \href{https://www.publichealth.columbia.edu/research/population-health-methods/path-analysis\#:~:text=Path\%20analysis\%20is\%20always\%20theory,remembered\%20when\%20performing\%20path\%20analysis}{\ul{https://www.publichealth.columbia.edu/research/population-health-methods/path-analysis}}.
\item
  Al-Hasi, D. (2023). \emph{Body Dysmorphic Disorder and Body Esteem and
  the Relationship to Self-Harm among a Sample of Patients in Cosmetic
  Clinics}. Unpublished Master Thesis, Al-Ahliyya Amman University.
\item
  Shawaqfeh, B. \& Almahaireh, A. (\RL{2022}). Social Anxiety Disorder
  Symptoms among Divorced Women in Amman and the Differences in It
  According to Some Demographic Variables. \emph{Al-Quds Open University
  Journal For Research}, 13 (38), 87-96.
\item
  Abu Al-wafaa, A. (2024). The Relationship Between Social Anxiety and
  Quality of Life Among Sample of Psychiatric Clinic Clients in Amman: a
  Comparative Correlation Study. Unpublished Master Thesis. The
  University of Jordan.
\item
  Garson, G. D. (2014). Path analysis (2\textsuperscript{nd} ed.).
  Asheboro, NC: Statistical Associates Publishing.
\item
  Hair, J. F., Black, W. C., Babin, B., Anderson, R. E., \& Tatham, R.
  (2018). \emph{Multivariate data analysis}: Cengage.
\item
  Suluhan, D., SAĞIROĞLU, E., Demir, S., YILDIZ, D., Onat, M., \& Şenel,
  E. (2023). Determination of the relationship between self-esteem and
  social anxiety in adolescents with burns. \emph{Türkiye Çocuk
  Hastalıkları Dergisi}, 1-7.
\item
  Gorbani, A., Rezaiee Moradali, M., \& Shabanloei, R. (2021).
  Relationship between self‐esteem and perceived social support in burn
  patients in Sina Hospital of Tabriz. \emph{Nursing open}, 8(3),
  1194-1200.
\item
  Hussain, M., Tariq, M., \& Hussain, M. (2023). Social appearance
  anxiety, psychological distress and quality of life among patients
  with burn injuries. \emph{Journal of Professional \& Applied
  Psychology}, 4(3), 418-428.
\item
  Yeşilyurt, İ., \& Kendirkıran, G. (2024). The effect of social
  appearance anxiety and body perception on the quality of life in burn
  patients. \emph{International Wound Journal}, 21(2), e14720.
\end{enumerate}

\textbf{\hfill\break
}

\textbf{Authors Biography}

{\def\LTcaptype{none} % do not increment counter
\begin{longtable}[]{@{}
  >{\centering\arraybackslash}p{(\linewidth - 2\tabcolsep) * \real{0.3832}}
  >{\raggedright\arraybackslash}p{(\linewidth - 2\tabcolsep) * \real{0.5716}}@{}}
\toprule\noalign{}
\begin{minipage}[b]{\linewidth}\centering
\includegraphics[width=0.88699in,height=0.93539in]{media/image9.jpeg}
\end{minipage} & \begin{minipage}[b]{\linewidth}\raggedright
Abdallah Almahaireh is Counselling and Special Education Department,
School of Educational Science, The University of Jordan. He received the
Ph.D. degree in Educational and Psychological Counselling at The
University of Jordan.
\end{minipage} \\
\midrule\noalign{}
\endhead
\bottomrule\noalign{}
\endlastfoot
\includegraphics[width=0.83296in,height=1.04167in]{media/image10.jpeg} &
Baha\textquotesingle{} Shawaqfeh is Full-time lecturer in Counselling
and Mental Health Department, Faculty of Educational Science, The World
Islamic Sciences and Education University. He received the Ph.D. degree
in Educational and Psychological Counselling at The University of
Jordan. \\
\includegraphics[width=1.01996in,height=1.01783in]{media/image11.jpeg} &
Bahaa Eldeen Abu Serdaneh has a Ph.D. degree in Educational Psychology,
School of Educational Science, at The University of Jordan. \\
\end{longtable}
}

\end{document}
